\documentclass{beamer} 
\usepackage{beamerthemesplit} 
\usepackage{wrapfig} 
\usepackage{verbatim} 
\usetheme{SPbGU} 
\usepackage{pdfpages} 
\usepackage{amsmath} 
\usepackage{cmap}
\usepackage{array} 
\usepackage[T2A]{fontenc} 
\usepackage[utf8]{inputenc} 
\usepackage{tikz} 
\usetikzlibrary{positioning, automata}
\usepackage{multirow} 
\usepackage[noend]{algpseudocode} 
\usepackage{algorithm} 
\usepackage{algorithmicx} 
\usetikzlibrary{shapes,arrows} 
\usepackage{fancyvrb} 
\usepackage{pgfplots} 
\usepackage{sidecap} 
\usepackage{soul}
\usepackage{xcolor}
\usepackage{tabu}
\usepackage{tikz}
\usetikzlibrary{calc}
\usepackage{zref-savepos}
\usepackage{colortbl}
\usepackage[normalem]{ulem}
\pgfplotsset{compat=1.9} 
\newtheorem{rutheorem}{Теорема} 
\newtheorem{ruproof}{Доказательство} 
\newtheorem{rudefinition}{Определение} 
\newtheorem{rulemma}{Лемма} 
\beamertemplatenavigationsymbolsempty 

\newcounter{NoTableEntry}
\renewcommand*{\theNoTableEntry}{NTE-\the\value{NoTableEntry}}

\newcommand*{\strike}[2]{%
	\multicolumn{1}{#1}{%
		\stepcounter{NoTableEntry}%
		\vadjust pre{\zsavepos{\theNoTableEntry t}}% top
		\vadjust{\zsavepos{\theNoTableEntry b}}% bottom
		\zsavepos{\theNoTableEntry l}% left
		\hspace{0pt plus 1filll}%
		#2% content
		\hspace{0pt plus 1filll}%
		\zsavepos{\theNoTableEntry r}% right
		\tikz[overlay]{%
			\draw
			let
			\n{llx}={\zposx{\theNoTableEntry l}sp-\zposx{\theNoTableEntry r}sp-\tabcolsep},
			\n{urx}={\tabcolsep},
			\n{lly}={\zposy{\theNoTableEntry b}sp-\zposy{\theNoTableEntry r}sp},
			\n{ury}={\zposy{\theNoTableEntry t}sp-\zposy{\theNoTableEntry r}sp}
			in
			(\n{llx}, \n{lly}) -- (\n{urx}, \n{ury})
			;
		}% 
	}%
}

\title[Effective CFL-reachability]{Эффективные алгоритмы для CFL-reachability} 
% То, что в квадратных скобках, отображается в левом нижнем углу. 
\institute[]{ JetBrains Research, Programming Languages and Tools Lab  \\
    Сенкт-Петербургский Академический Университет } 

\author[Екатерина Шеметова]{Екатерина Шеметова}

\date{14.12.2019} 

\begin{document} 
\definecolor{red}{RGB}{255,0,0} 

\begin{frame}
	\begin{tabular}{p{2.0cm} p{7.5cm} p{1cm}}
        \begin{center}
            \includegraphics[height=1.5cm]{pictures/jetbrainsResearch.pdf}
        \end{center}
        
    \end{tabular}
	\titlepage
\end{frame}


\begin{frame}
\frametitle{CFL-reachability}
        \begin{center}
           достижимость в графе + контекстно-свободные ограничения
        \end{center}
\begin{columns}[c] % The "c" option specifies centered vertical alignment while the "t" option is used for top vertical alignment

\column{.55\textwidth} % Left column and width
\begin{tikzpicture}[shorten >=1pt,node distance=2cm,on grid,auto] 
   \node[state] (q_1)   {$1$}; 
   \node[state] (q_2) [above right=of q_1] {$2$}; 
   \node[state] (q_3) [right=of q_2] {$3$}; 
   \node[state] (q_4) [right=of q_3] {$4$};
    \path[->] 
    (q_1) edge  node {a} (q_2)          
    (q_2) edge  node {a} (q_3)
    (q_3) edge  node {a} (q_1)
    (q_3) edge[bend left, above]  node {b} (q_4)
    (q_4) edge[bend left, below]  node {b} (q_3);
\end{tikzpicture}

\column{.4\textwidth} % Right column and width
$S \rightarrow aSb$ $|$ $ab$
\\
Решения: $(S, 2, 4), (S, 1, 3), ..., (S, 3, 3)$
\\2-4: ab
\\1-3: aabb
\\3-4: aaabbb
\\ ...
\\3-3: aaaaaabbbbbb


\end{columns}


\begin{itemize}
\item Вход: контекстно-свободная грамматика $G$ и помеченный ориентированный граф $D$
\itemСферы применения: статический анализ кода, запросы к графовым базам данных
\end{itemize}
\end{frame}

\begin{frame}
\frametitle{Проблемы и цель}
\textbf{Цель:} Хотим эффективно решать задачу CFL-reachability
\\
\textbf{Проблемы:} 
\begin{enumerate}
\item Классический алгоритм для решения задачи CFL-reachability работает за $O(n^3)$ ($n$ --- число вершин в графе).
\\\textbf{Большая открытая проблема:} существует ли алгоритм, работающий за $O(n^{3-\varepsilon})$ (т.н. \textit{субкубический алгоритм})?

\end{enumerate}
\end{frame}

\begin{frame}
\frametitle{Проблемы и цель}
\textbf{Цель:} Хотим эффективно решать задачу CFL-reachability
\\ 
\textbf{Проблемы:} 
\begin{enumerate}
\item Классический алгоритм для решения задачи CFL-reachability работает за $O(n^3)$ ($n$ --- число вершин в графе).
\\\textbf{Большая открытая проблема:} существует ли алгоритм, работающий за $O(n^{3-\varepsilon})$ (т.н. \textit{субкубический алгоритм})?
\item Параллельный алгоритм --- доказано, что для задачи CFL-reachability нельзя построить эффективный параллельный алгоритм

\end{enumerate}
\end{frame}

\begin{frame}
\frametitle{Проблемы и цель}
\textbf{Цель:} Хотим эффективно решать задачу CFL-reachability
\\ 
\textbf{Проблемы:} 
\begin{enumerate}
\item Классический алгоритм для решения задачи CFL-reachability работает за $O(n^3)$ ($n$ --- число вершин в графе).
\\\textbf{Большая открытая проблема:} существует ли алгоритм, работающий за $O(n^{3-\varepsilon})$ (т.н. \textit{субкубический алгоритм})?
\item Параллельный алгоритм --- доказано, что для задачи CFL-reachability нельзя построить эффективный параллельный алгоритм

\end{enumerate}
\textbf{Подход:} Попробуем найти подклассы граматик, для которых задача может быть эффективно решена



\end{frame}

\begin{frame}
\frametitle{Параллельная сложность}
\begin{itemize}
\item Хотим знать, какие задачи эффективно параллелятся, а какие --- нет
\item Что в теории значит ``эффективно параллелялятся''?
\end{itemize}
\end{frame}
\begin{frame}
\frametitle{Параллельная сложность}
\begin{itemize}
\item Хотим знать, какие задачи эффективно параллелятся, а какие --- нет
\item Что в теории значит ``эффективно параллелялятся''?
\begin{itemize}
\item Для входа длиной $n$ задача может быть решена:
\begin{itemize}
\item полиномиальным от $n$ числом процессоров и 
\item за полилогарифмическое от $n$ время $polylog(n)$
\end{itemize}
\item Если время хотя бы линейно от входа, то параллелить уже неэффективно
\end{itemize}
\end{itemize}
\end{frame}

\begin{frame}
\frametitle{CFL-reachability и параллельная сложность}
\begin{itemize}
\item Классический результат: в общем случае для задачи CFL-reachability нельзя построить эффективный параллельный алгоритм
\item Но утверждается, что для некоторых фиксированных классов контекстно-свободных грамматик данная задача эффективно параллелится
\end{itemize}
\end{frame}


\begin{frame}
\frametitle{CFL-reachability и параллельная сложность}
\begin{itemize}
\item Классический результат: в общем случае для задачи CFL-reachability нельзя построить эффективный параллельный алгоритм
\item Но утверждается, что для некоторых фиксированных классов контекстно-свободных грамматик данная задача эффективно параллелится
\item \textbf{Вопросы:} Что объединяет эти эффективные классы? Много ли их? Почему они обладают этим свойством? Можем ли мы получить новые, полезные на практике?
\end{itemize}
\end{frame}




\begin{frame}
\frametitle{Эффективные подклассы}
\begin{itemize}
\item Эффективность зависит от специального параметра языка --- рационального индекса $\rho_L(n)$
\item $\rho_L(n) = \max\{\min\{|w|:w \in L \cap K\}, K \in {Rat}_n, L \cap K \neq \emptyset\}$
\begin{columns}[c] % The "c" option specifies centered vertical alignment while the "t" option is used for top vertical alignment

\column{.55\textwidth} % Left column and width
\begin{tikzpicture}[shorten >=1pt,node distance=2cm,on grid,auto] 
   \node[state] (q_1)   {$1$}; 
   \node[state] (q_2) [above right=of q_1] {$2$}; 
   \node[state] (q_3) [right=of q_2] {$3$}; 
   \node[state] (q_4) [right=of q_3] {$4$};
    \path[->] 
    (q_1) edge  node {a} (q_2)          
    (q_2) edge  node {a} (q_3)
    (q_3) edge  node {a} (q_1)
    (q_3) edge[bend left, above]  node {b} (q_4)
    (q_4) edge[bend left, below]  node {b} (q_3);
\end{tikzpicture}

\column{.4\textwidth} % Right column and width
$S \rightarrow aSb$ $|$ $ab$
\\
Решения: $(S, 2, 4), (S, 1, 3), ..., (S, 3, 3)$
\\2-4: ab
\\1-3: aabb
\\3-4: aaabbb
\\ ...
\\\textbf{3-3: aaaaaabbbbbb}


\end{columns}
\end{itemize}
\end{frame}

\begin{frame}
\frametitle{Эффективные подклассы}
\begin{itemize}
\item Эффективность зависит от специального параметра языка --- рационального индекса  $\rho_L(n)$
\item $\rho_L(n) = \max\{\min\{|w|:w \in L \cap K\}, K \in {Rat}_n, L \cap K \neq \emptyset\}$
\item Если индекс полиномиален от числа вершин в графе --- задача может быть эффективно распараллелена, экспоненциален --- нет
\item Что влияет на величину индекса?
\end{itemize}
\end{frame}


\begin{frame}
\frametitle{Эффективные подклассы}
\begin{itemize}
\item Эффективность зависит от специального параметра языка --- рационального индекса  $\rho_L(n)$
\item $\rho_L(n) = \max\{\min\{|w|:w \in L \cap K\}, K \in {Rat}_n, L \cap K \neq \emptyset\}$
\item Если индекс полиномиален от числа вершин в графе --- задача может быть эффективно распараллелена, экспоненциален --- уже нет
\item Что влияет на величину индекса?
\begin{itemize}
\item Некоторые ограничения на стек pushdown автомата делают рациональный индекс соответствующего языка полиномиальным!
\item Примеры: one-turn, oscillation-boundness, один элемент в стековом алфавите и т.д.
\end{itemize}
\item Также полезно изучить влияние операций над языками на величину индекса
\end{itemize}
\end{frame}

\begin{frame}
\frametitle{Результаты}
\begin{itemize}
\item Ограничения на стек  pushdown автомата дают нам эффективные языки
\item Найдены новые подклассы, для которых можно построить эффективный параллельный алгоритм: металинейные, суперлинейные, oscillation-bounded языки
\item А также известные эффективные языки можно комбинировать с помощью операций конкатенации, union, пересечения с регулярными языками, substitution closure и др. и получать новые эффективные примеры
\item Для решения задачи адаптирован параллельный алгоритм Брента--Гольдшляггера--Риттера  (на эффективных подклассах число процессоров $O(n^6)$, время --- $O(\log^2 n)$, $n$ --- число вершин в графе)
\end{itemize}
\end{frame}

\begin{frame}
\frametitle{Результаты}
А что по поводу субкубического алгоритма?
\begin{itemize}
\item Эффективные для распараллеливания подклассы языков более эффективны и для последовательного алгоритма
\item Благодаря использованию класса one-counter языков, была получена аппроксимация общей задачи за субкубическое время
\end{itemize}
\end{frame}

\begin{frame}
\frametitle{Публикации}
\begin{itemize}
\item \textbf{Журнал ``Theory of Computing Systems''} \\
Ekaterina Shemetova, Semyon Grigorev and Alexander Okhotin. Parallel efficient variants of context-free language reachability problem.
\\\textit{Статус: планируется публикация}
\end{itemize}
Прошлое полугодие: 
\begin{itemize}
\item  \textbf{Журнал ``Труды Института системного программирования РАН''}
\\Шеметова Е.Н., Григорьев С.В. Задача поиска путей в ациклических графах с ограничениями в терминах булевых грамматик. 
\\\textit{Статус: опубликовано}
\end{itemize}

\end{frame}
\begin{frame}
\frametitle{Планы}
\begin{itemize}
\item Построение обощенного субкубического алгоритма для найденных эффективных подклассов КС-языков
\item Планируется найти эффективные подклассы графов путем изучения влияния структурных характеристик графа, таких как treewidth, pathwidth, сutwidth и др., на эффективность решения задачи
\item Пробуем использовать найденные эффективные подклассы (как графов, так и грамматик) для получения решения задачи в общем случае 
\item Подобных проблем много (например APSP), и занимается ими целое сообщество (``fine-grained'' complexity). Ждем ответа от Barna Saha (University of California Berkeley, автор cубкубического алгоритма для LED).
\end{itemize}

\end{frame}
\end{document}

